% STATUS: written 2026-11-05 (Freudenburg 2017, Drensky--Makar-Limanov 2009,
% Drensky--Gupta 2005, Baltazar 2016, Mendes--Pan 2017, Rewri--Kour 2025,
% Moulin Ollagnier--Nowicki 2011 and Bedratyuk 2005 read in full at the cited pages),
% with additions of 2026-11-06 from Mishra--Mondal--Shukla 2026 (pp. 5, 11, 14).
% Sources: Freudenburg2017, DrenskyMakarLimanov2009, DrenskyGupta2005,
% Baltazar2016, MendesPan2017, RewriKour2025, MoulinOllagnierNowicki2011,
% Bedratyuk2005Casimir, MishraMondalShukla2026Monomial, Li2015DarbouxMonomial,
% MishraMondalShukla2025Simple.
% Origins: the named families are defined from Freudenburg 2017 and the papers
% that carry them; Weitzenbock 1932 and Shamsuddin 1977 could not be downloaded
% and are cited through the recent sources that quote them (recorded below).
\section{Derivations of polynomial algebras}\label{sec:polynomial-weyl:polynomial}

The chapter returns to the polynomial ring, the first home of the
derivative. The fourteenth problem of Hilbert asked whether certain rings of
invariants are finitely generated. The Jacobian conjecture asked whether a
polynomial map with a unit Jacobian determinant is invertible
\cite[p.~xii]{Freudenburg2017}. Derivations answer on both fronts, and
the answers organize the section: kernels, the named families, and the
kernel theorems.

The named families live on $k[x_{1},\dots,x_{n}]$ with $k$ a field of
characteristic zero. Section~\ref{sec:basic-theory:exponential} met the
locally nilpotent derivation in general; Freudenburg makes it concrete on
polynomial rings \cite[p.~76]{Freudenburg2017}. Recall: a derivation $D$
is locally nilpotent when every element is killed by a high enough power
of $D$ \cite[p.~21]{Freudenburg2017}. % claim-ok: recall of the source page of Definition def:basic-nilpotent, Freudenburg 2017 principle 2 region

\begin{definition}[{\cite[p.~76]{Freudenburg2017}}]\label{def:pw:triangular}
Let $B=k[x_{1},\dots,x_{n}]$ and $D\in\mathrm{Der}_{k}(B)$. The derivation
$D$ is a \emph{triangular derivation} of $B$ if $D(x_{1})\in k$ and
$D(x_{i})\in k[x_{1},\dots,x_{i-1}]$ for $i=2,\dots,n$. The derivation $D$
is a \emph{monomial derivation} if each image $D(x_{i})$ is a monomial in
$x_{1},\dots,x_{n}$. The derivation $D$ is an \emph{elementary derivation}
if for some $j$ with $1\leq j\leq n$, $D(x_{i})=0$ for $1\leq i\leq j$ and
$D(x_{i})\in k[x_{1},\dots,x_{j}]$ for $j+1\leq i\leq n$.
\end{definition}

Triangularity depends on the choice of coordinates, and every triangular,
elementary or nice derivation of $k[x_{1},\dots,x_{n}]$ is locally
nilpotent \cite[p.~76]{Freudenburg2017}. Every elementary derivation is
both triangular and nice \cite[p.~76]{Freudenburg2017}. In one direction
the converse holds after a permutation of the variables, as
Proposition~\ref{prop:pw:monomial-triangular} states
\cite[p.~76]{Freudenburg2017}.

\begin{proposition}[{\cite[Proposition~3.4, p.~76]{Freudenburg2017}}]\label{prop:pw:monomial-triangular}
For the polynomial ring $B=k[x_{1},\dots,x_{n}]$, every locally nilpotent
monomial derivation of $B$ is triangular relative to some ordering of
$x_{1},\dots,x_{n}$.
\end{proposition}

Monomial derivations have a second home: the body of constants.
Definition~\ref{def:pw:triangular} states them through the images of
single variables, and Moulin Ollagnier and Nowicki study them for their
rings of constants \cite[p.~3138]{MoulinOllagnierNowicki2011}. Their
setting adds the field of rational constants, and a classical example of
the family sits inside it \cite[p.~3139]{MoulinOllagnierNowicki2011}.

\begin{example}[{\cite[p.~3139]{MoulinOllagnierNowicki2011}}]\label{ex:pw:jouanolou}
The \emph{Jouanolou derivation} with integer parameters $n\geq 2$ and
$s\geq 1$ is the derivation $J(n,s)$ of $k[x_{1},\dots,x_{n}]$ given by
$d(x_{1})=x_{2}^{s}$, $d(x_{2})=x_{3}^{s}$, \dots, $d(x_{n-1})=x_{n}^{s}$
and $d(x_{n})=x_{1}^{s}$. Jouanolou proved in 1979 that $J(3,s)$ has no
nontrivial Darboux polynomial for every $s\geq 2$
\cite[p.~3139]{MoulinOllagnierNowicki2011}.
\end{example}

Li carries the question of Darboux polynomials to a monomial derivation of
four variables \cite[p.~1]{Li2015DarbouxMonomial}. A polynomial $F$ is a
Darboux polynomial of $d$ when $d(F)=\Lambda F$ for some polynomial
$\Lambda$ \cite[p.~1]{Li2015DarbouxMonomial}. For a normal monomial
derivation with $\det(\beta_{ij}-I)\neq 0$, the absence of Darboux
polynomials is the same as a trivial field of constants
\cite[p.~2]{Li2015DarbouxMonomial}. Li treats the case where that
determinant vanishes, for the derivation
$d(x)=t^{2}$, $d(y)=zt$, $d(z)=y^{2}$, $d(t)=xy$
\cite[p.~2]{Li2015DarbouxMonomial}.

\begin{theorem}[{\cite[Theorem~2.3, p.~3]{Li2015DarbouxMonomial}}]\label{thm:pw:darboux-fours}
Let $d$ be the monomial derivation of $k[x,y,z,t]$ with $d(x)=t^{2}$,
$d(y)=zt$, $d(z)=y^{2}$ and $d(t)=xy$. Then $d$ has no Darboux polynomial
if and only if its field of constants is trivial.
\end{theorem}

So the equivalence survives the vanishing of the determinant, for this one
case. The proof restricts a possible Darboux polynomial to its homogeneous
parts and finds a contradiction from the degrees
\cite[p.~3]{Li2015DarbouxMonomial}.

The nicest monomial way into the kernel is linear. A derivation of
$k[x_{1},\dots,x_{n}]$ is linear when it maps the variables into their own
linear span \cite[p.~394]{DrenskyGupta2005}, and a Weitzenb\"ock
derivation is the linear locally nilpotent case. Weitzenb\"ock studied it
in 1932, and the name comes with his theorem
\cite[p.~41]{DrenskyMakarLimanov2009}. His paper could not be downloaded
for this book; the definition and the theorem are quoted here from Drensky
and Makar-Limanov.

\begin{definition}[{\cite[p.~41]{DrenskyMakarLimanov2009}}]\label{def:pw:weitzenbock}
A derivation $\Delta$ of the polynomial algebra $K[Z]=K[z_{1},\dots,z_{m}]$
over a field $K$ of characteristic zero is a \emph{Weitzenb\"ock
derivation} if it is linear and locally nilpotent, that is, $\Delta$ is a
nilpotent linear operator on the vector space
$Kz_{1}\oplus\cdots\oplus Kz_{m}$ extended to a derivation of $K[Z]$.
\end{definition}

The next theorem is the finiteness result of Weitzenb\"ock.

\begin{theorem}[{\cite[Weitzenb\"ock 1932, p.~42]{DrenskyMakarLimanov2009}}]\label{thm:pw:weitzenbock}
If $\Delta$ is a Weitzenb\"ock derivation of $K[Z]$, then the algebra of
constants $K[Z]^{\Delta}$ is a finitely generated algebra.
\end{theorem}

A Weitzenb\"ock derivation is determined up to a change of basis by its
Jordan normal form, whose cells are upper triangular with zero diagonal
\cite[p.~42]{DrenskyMakarLimanov2009}. For two-by-two cells only, Nowicki
made a conjecture on the generators of the kernel, and Drensky and
Makar-Limanov proved it \cite[pp.~42--45]{DrenskyMakarLimanov2009}.

\begin{theorem}[{\cite[Theorem~3, p.~44]{DrenskyMakarLimanov2009}}]\label{thm:pw:nowicki}
Let $X=\{x_{1},\dots,x_{n}\}$, $Y=\{y_{1},\dots,y_{n}\}$, $n\geq 1$, and
let $\Delta$ be the Weitzenb\"ock derivation of $K[X,Y]$ defined by
$\Delta=\sum_{i=1}^{n} x_{i}\,\partial/\partial y_{i}$. The algebra of
constants $K[X,Y]^{\Delta}$ is generated by $x_{i}$, $i=1,\dots,n$, and
$u_{ij}=x_{i}y_{j}-x_{j}y_{i}$, $1\leq i<j\leq n$.
\end{theorem}

The generators of such a kernel can also be built one at a time. Bedratyuk
copies the Casimir elements of a Lie algebra over to a derivation
\cite[p.~2]{Bedratyuk2005Casimir}. For a derivation $D$ of $k[X]$ he takes
two dual $D$-modules $U$ and $U^{*}$, with bases $u_{i}$ and $u^{*}_{i}$, and
forms the element $\Delta(U,U^{*})=\sum_{i}u_{i}u^{*}_{i}$
\cite[Definition~2.4, p.~2]{Bedratyuk2005Casimir}. That element lies in the
kernel of $D$ \cite[Theorem~2.2, p.~2]{Bedratyuk2005Casimir}.

\begin{theorem}[{\cite[Theorem~2.2, p.~2]{Bedratyuk2005Casimir}}]\label{thm:pw:casimir-in-kernel}
Let $D$ be a derivation of $k[X]$ and let $U$ and $U^{*}$ be two dual
$D$-modules in $k[X]$. Then the Casimir element $\Delta(U,U^{*})$ belongs to
the kernel $k[X]^{D}$.
\end{theorem}

For a linear derivation the converse holds, and every homogeneous element of
the kernel is such an element
\cite[Theorem~2.3, pp.~2--3]{Bedratyuk2005Casimir}. The Weitzenb\"ock
derivation of Definition~\ref{def:pw:weitzenbock} is linear, so its kernel is
generated by Casimir elements \cite[pp.~3--4]{Bedratyuk2005Casimir}. That
gives an algorithm, and Bedratyuk runs it to list the kernel for at most six
variables \cite[p.~1]{Bedratyuk2005Casimir}.

The kernel theorems of Weitzenb\"ock and Nowicki describe the constants of
a linear derivation. Nonlinear families raise a second question: which
derivations of $k[x,y]$ generate no invariant ideals at all
\cite[p.~2031]{Baltazar2016}? A derivation $d$ of a ring is simple when
the only $d$-stable ideals are $0$ and the ring itself
\cite[p.~2031]{Baltazar2016}; The class is named after Shamsuddin, and Lequain characterizes
which members of it are simple derivations
\cite[pp.~801--807]{Lequain2008Simple}.

\begin{definition}[{\cite[p.~2033]{Baltazar2016}}]\label{def:pw:shamsuddin}
A derivation $d$ of $k[x,y]$ is a \emph{Shamsuddin derivation} if it has
the form
\begin{equation*}
d=\partial_{x}+(a(x)\,y+b(x))\,\partial_{y},
\end{equation*}
where $a(x),b(x)\in k[x]$.
\end{definition}

Among all derivations of $k[x,y]$, the simple ones turn out to be lonely:
Mendes and Pan prove that no nontrivial automorphism of $k[x,y]$ commutes
with them. Write $\mathrm{Aut}(D)$ for the subgroup of
$\mathrm{Aut}_{K}(K[x,y])$ whose elements commute with $D$; this is the
\emph{isotropy group} of $D$ \cite[p.~2]{MendesPan2017}. Their theorem
reads as follows \cite[p.~2]{MendesPan2017}.

\begin{theorem}[{\cite[Theorem~1, p.~2]{MendesPan2017}}]\label{thm:pw:isotropy-simple}
Let $K$ be an algebraically closed field of characteristic zero and let
$D$ be a simple derivation of $K[x,y]$. Then the isotropy group
$\mathrm{Aut}(D)$ of $D$ in $\mathrm{Aut}_{K}(K[x,y])$ is trivial.
\end{theorem}

For Shamsuddin derivations the converse holds under a hypothesis, as the
companion theorem states \cite[p.~2]{MendesPan2017}.

\begin{theorem}[{\cite[Theorem~2, p.~2]{MendesPan2017}}]\label{thm:pw:shams-simple}
Let $D$ be a Shamsuddin derivation $D=\partial_{x}+(ay+b)\partial_{y}$ of
$K[x,y]$ with $a\neq 0$. Then $D$ is simple if and only if
$\mathrm{Aut}(D)=\{\mathrm{id}\}$.
\end{theorem}

A third named family comes from ecology \cite[p.~2]{RewriKour2025}. It
looks quadratic rather than linear, and it models the predator and prey
populations through a system of differential equations
\cite[p.~2]{RewriKour2025}.

\begin{definition}[{\cite[(1), p.~2]{RewriKour2025}}]\label{def:pw:lotka-volterra}
Let $K$ be a field of characteristic zero, $n\geq 3$, and let the indices
of the variables $x_{1},\dots,x_{n}$ of $K[X]=K[x_{1},\dots,x_{n}]$ be
regarded as elements of $\mathbb{Z}_{n}$. A $K$-derivation
$d\colon K[X]\to K[X]$ of the form
\begin{equation*}
d(x_{i})=x_{i}\,(x_{i-1}-C_{i}\,x_{i+1})
\end{equation*}
where $C_{i}\in K$ for all $i\in\mathbb{Z}_{n}$, is a
\emph{Lotka--Volterra derivation}.
\end{definition}

Lotka--Volterra derivations are not simple \cite[p.~2]{RewriKour2025}, so
Theorem~\ref{thm:pw:isotropy-simple} does not speak for them. Rewri and
Kour compute their isotropy groups directly
\cite[pp.~21--23]{RewriKour2025}.

\begin{theorem}[{\cite[Theorem~4.1, p.~21]{RewriKour2025}}]\label{thm:pw:lv-finite}
Let $d$ be the Lotka--Volterra derivation of $K[X]$ as in
Definition~\ref{def:pw:lotka-volterra}. If $C_{i}\in K^{*}$ for all $i$,
then the isotropy group of $d$ is finite.
\end{theorem}

When all the coefficients agree and lie in $\{\pm1\}$, the group is dihedral.

\begin{corollary}[{\cite[Corollary~4.3, p.~23]{RewriKour2025}}]\label{cor:pw:lv-dihedral}
Let $d$ be the Lotka--Volterra derivation of $K[X]$ as in
Definition~\ref{def:pw:lotka-volterra}. If $C_{1}=\dots=C_{n}$ and
$C_{1}\in\{\pm 1\}$, then the isotropy group of $d$ is isomorphic to the
dihedral group of order $2n$.
\end{corollary}

The monomial family of Definition~\ref{def:pw:triangular} also has a known
isotropy group. Mishra, Mondal and Shukla treat a monomial derivation whose
exponent matrix has nonzero determinant
\cite[Theorem~3.1, p.~5]{MishraMondalShukla2026Monomial}. Such an
automorphism acts on the variables one at a time, up to a diagonal change
\cite[Theorem~3.1, p.~5]{MishraMondalShukla2026Monomial}.

\begin{theorem}[{\cite[Theorem~3.1, p.~5]{MishraMondalShukla2026Monomial}}]\label{thm:pw:monomial-isotropy}
Let $k$ be a field of characteristic zero and let $d$ be the monomial
derivation of $k[x_{1},\dots,x_{n}]$ with $d(x_{i})=\alpha_{i}x_{1}^{\beta_{i1}}\cdots x_{n}^{\beta_{in}}$
for $\alpha_{i}\in k^{*}$ and non-negative integers $\beta_{ij}$. Put
$A=(\beta_{ij}-\delta_{ij})$ and suppose that $\beta_{ii}\geq 1$ for all
$i$, that $\det A\neq 0$, and that the field of rational constants of $d$ is
$k$. Let $\rho\in\operatorname{Aut}(k[X])$. Then $\rho$ commutes with $d$ if
and only if there are a permutation $\sigma$ and constants $c_{1},\dots,c_{n}\in k^{*}$
with $\rho(x_{i})=c_{i}x_{\sigma(i)}$, $P_{\sigma}A=AP_{\sigma}$ and
$\prod_{t}c_{t}^{a_{ti}}=\alpha_{\sigma(i)}/\alpha_{i}$ for all $i$.
\end{theorem}

For the Jouanolou-type derivation $J=\sum_{i}x_{i}^{r}\partial_{x_{\sigma(i)}}$
the commuting automorphisms are again monomial
\cite[Theorem~4.2, p.~11]{MishraMondalShukla2026Monomial}. Each one is fixed
by a permutation that commutes with $\sigma$ and by constants that satisfy
$\lambda_{\sigma(i)}=\lambda_{i}^{r}$ \cite[Theorem~4.2, p.~11]{MishraMondalShukla2026Monomial}.
The full group is then a semidirect product written out cycle by cycle
\cite[Theorem~4.6, p.~14]{MishraMondalShukla2026Monomial}.

A simple derivation that takes no unit value has a small group too.
Mishra, Mondal and Shukla extend a non-invertible simple derivation of an
initial block of variables to more variables
\cite[p.~1]{MishraMondalShukla2025Simple}. Under their conditions the
commuting group of the extension lies in the translations
\cite[Theorem~2.4, p.~5]{MishraMondalShukla2025Simple}.

\begin{theorem}[{\cite[Theorem~2.4, p.~5]{MishraMondalShukla2025Simple}}]\label{thm:pw:noninvertible-simple}
Let $i,n$ be positive integers with $i\geq 2$ and $n>i$. Let $d_{i}$ be a
simple derivation of $k[x_{1},\dots,x_{i}]$ that satisfies the three
conditions of the cited theorem, and let $d_{n}$ extend it to
$k[x_{1},\dots,x_{n}]$ with $d_{n}(x_{j})\in k[x_{j-1}]\setminus k$ for
$i+1\leq j\leq n$. Then the commuting group of $d_{n}$ is
$\{(x_{1},\dots,x_{n-1},x_{n}+c)\mid c\in k\}$.
\end{theorem}

So the isotropy group is as small as it can be, a one-parameter group of
translations \cite[p.~1]{MishraMondalShukla2025Simple}. This continues the
pattern of Theorem~\ref{thm:pw:isotropy-simple}: simplicity leaves almost no
symmetry.

The named families of the section return again in the geometry of
polynomial maps. Rentschler classified the planar case in 1968: every
nonzero locally nilpotent derivation of $k[x,y]$ is conjugate, by a tame
automorphism, to a triangular one \cite[Theorem~4.1, p.~113]{Freudenburg2017}.

\begin{theorem}[{\cite[Theorem~4.1, p.~113]{Freudenburg2017}}]\label{thm:pw:rentschler}
Let $k[x,y]=k^{[2]}$ and let $D\in\mathrm{LND}(k[x,y])$ be nonzero. There
exists $p(x)\in k[x]$ and a tame automorphism $\alpha\in\mathrm{GA}_{2}(k)$
such that $\alpha D\alpha^{-1}=p(x)\,\partial_{y}$.
\end{theorem}

the theorem of Rentschler also gives the theorem of Jung on planar automorphisms as
a corollary \cite[p.~113]{Freudenburg2017}. The Jacobian conjecture closes
the section: it asks whether a polynomial endomorphism $F$ of
$\mathbb{A}^{n}_{k}$ with $\det J(F)$ a nonzero constant is an automorphism
\cite[p.~1]{BelovKanelKontsevich2007}. The conjecture has a formulation in
derivations: if the Jacobian derivation $\Delta_{f}$ of
$f=(f_{1},\dots,f_{n-1})$ has a slice $s$, then $k[f_{1},\dots,f_{n-1},s]$
is the whole polynomial ring \cite[p.~81]{Freudenburg2017}. The section
that follows moves the same question to the Weyl algebra, where it is the
Dixmier conjecture \cite[p.~241]{Dixmier1968}.

Dixmier's problem 6 is settled for somewhat commutative algebras and problem 3 for the ring of differential operators, which are the two questions that chapter records about a derivation algebra being determined by its invariants~\cite{WlDixmierSProblem6ForSomewhatCom2005}.

Closed polynomials are used to compute the kernel of a monomial derivation, which is the concrete form of the kernel question the chapter records for polynomial algebras~\cite{WlClosedPolynomialsAndTheirAppli2019}.

The images of higher-order differential operators of a polynomial algebra are determined, and in positive characteristic the image of the differential operators is shown to have a stated form~\cite{WlImagesOfHigherOrderDifferentia2019}.

For an algebraically closed field $\mathbb{K}$ of characteristic zero and the Lie algebra $W_n$ of all $\mathbb{K}$-derivations of the polynomial ring in $n$ variables over $\mathbb{K}$, Makedonskyi proves that every Lie algebra of dimension $n$ over $\mathbb{K}$ embeds in $W_n$ so that any basis of its image is a basis of the free $W_n$-module~\cite{WlOnNoncommutativeBasesOfFreeMod2015}.

Makedonskyi proves that over an algebraically closed field of characteristic zero every Lie algebra of dimension $n$ embeds in the Lie algebra $W_n$ of all derivations of the polynomial ring in $n$ variables in such a way that any basis of its image is a basis of the free module $W_n$~\cite{WlOnNoncommutativeBasesOfFreeMod2015}.

Van den Essen, Wright and Zhao show that the image of any locally finite $k$-derivation of the polynomial algebra $k[x,y]$ in two variables over a field $k$ of characteristic zero is a Mathieu subspace, and that the two-dimensional Jacobian conjecture is equivalent to the corresponding statement for every $k$-derivation $D$ of $k[x,y]$ with $1$ in its image and zero divergence~\cite{WlImagesOfLocallyFiniteDerivatio2011}.

Jia, Du and Tian prove the $E$-derivation analogue of the result of van den Essen, Wright and Zhao: the image of a locally finite $K$-$E$-derivation of the polynomial algebra in two variables over a field $K$ of characteristic zero is a Mathieu subspace, and together with that result this confirms the LFED conjecture in the stated case~\cite{WlImagesOfLocallyFiniteEDerivati2023}.

Makedonskyi proves that over an algebraically closed field of characteristic zero every Lie algebra of dimension $n$ embeds in the Lie algebra $W_n$ of all derivations of the polynomial ring in $n$ variables in such a way that any basis of its image is a basis of the free module $W_n$~\cite{WlOnNoncommutativeBasesOfFreeMod2015}.
